\documentclass[10pt]{article}

\usepackage[margin=0.82in]{geometry}
\usepackage[T1]{fontenc}
\usepackage{lmodern}
\usepackage{microtype}
\usepackage{amsmath,amssymb,amsthm,mathtools,bm}
\usepackage{aliascnt}
\usepackage{booktabs,array,tabularx,longtable}
\usepackage{enumitem}
\usepackage{graphicx}
\usepackage{tikz}
\usetikzlibrary{arrows.meta,positioning,fit,calc}
\usepackage{xcolor}
\usepackage{natbib}
\usepackage[hidelinks]{hyperref}
\usepackage[nameinlink,noabbrev]{cleveref}

\setlength{\parindent}{0pt}
\setlength{\parskip}{4.2pt}
\setlist[itemize]{leftmargin=1.5em,itemsep=1.6pt,topsep=2pt}
\setlist[enumerate]{leftmargin=1.8em,itemsep=1.6pt,topsep=2pt}
\renewcommand{\arraystretch}{1.13}

\newtheorem{theorem}{Theorem}[section]

\newaliascnt{proposition}{theorem}
\newtheorem{proposition}[proposition]{Proposition}
\aliascntresetthe{proposition}

\newaliascnt{lemma}{theorem}
\newtheorem{lemma}[lemma]{Lemma}
\aliascntresetthe{lemma}

\newaliascnt{corollary}{theorem}
\newtheorem{corollary}[corollary]{Corollary}
\aliascntresetthe{corollary}

\newaliascnt{definition}{theorem}
\newtheorem{definition}[definition]{Definition}
\aliascntresetthe{definition}

\newaliascnt{assumption}{theorem}
\newtheorem{assumption}[assumption]{Assumption}
\aliascntresetthe{assumption}

\newaliascnt{example}{theorem}
\newtheorem{example}[example]{Example}
\aliascntresetthe{example}

\newaliascnt{remark}{theorem}
\newtheorem{remark}[remark]{Remark}
\aliascntresetthe{remark}

\crefname{assumption}{assumption}{assumptions}
\Crefname{assumption}{Assumption}{Assumptions}
\crefname{lemma}{lemma}{lemmas}
\Crefname{lemma}{Lemma}{Lemmas}
\crefname{proposition}{proposition}{propositions}
\Crefname{proposition}{Proposition}{Propositions}
\crefname{corollary}{corollary}{corollaries}
\Crefname{corollary}{Corollary}{Corollaries}
\crefname{theorem}{theorem}{theorems}
\Crefname{theorem}{Theorem}{Theorems}
\crefname{definition}{definition}{definitions}
\Crefname{definition}{Definition}{Definitions}
\crefname{example}{example}{examples}
\Crefname{example}{Example}{Examples}
\crefname{remark}{remark}{remarks}
\Crefname{remark}{Remark}{Remarks}

\newcommand{\R}{\mathbb{R}}
\newcommand{\N}{\mathbb{N}}
\newcommand{\E}{\mathbb{E}}
\newcommand{\Pp}{\mathbb{P}}
\newcommand{\cA}{\mathcal{A}}
\newcommand{\cC}{\mathcal{C}}
\newcommand{\cF}{\mathcal{F}}
\newcommand{\cH}{\mathcal{H}}
\newcommand{\cI}{\mathcal{I}}
\newcommand{\cM}{\mathcal{M}}
\newcommand{\cP}{\mathcal{P}}
\newcommand{\cR}{\mathcal{R}}
\newcommand{\cS}{\mathcal{S}}
\newcommand{\cZ}{\mathcal{Z}}
\newcommand{\Id}{\operatorname{Id}}
\newcommand{\im}{\operatorname{im}}
\newcommand{\Agg}{\operatorname{Agg}}
\newcommand{\Attn}{\operatorname{Attn}}
\newcommand{\FFN}{\operatorname{FFN}}
\newcommand{\GFFN}{\operatorname{GFFN}}
\newcommand{\Norm}{\operatorname{Norm}}
\newcommand{\Route}{\operatorname{Route}}
\newcommand{\softmax}{\operatorname{softmax}}
\newcommand{\argmin}{\operatorname*{arg\,min}}
\newcommand{\diam}{\operatorname{diam}}
\newcommand{\rank}{\operatorname{rank}}
\newcommand{\supp}{\operatorname{supp}}
\newcommand{\Span}{\operatorname{span}}
\newcommand{\one}{\mathbf{1}}
\newcommand{\given}{\,\middle|\,}
\newcommand{\norm}[1]{\left\lVert #1\right\rVert}
\newcommand{\abs}[1]{\left|#1\right|}
\newcommand{\ip}[2]{\left\langle #1,#2\right\rangle}
\newcommand{\fib}{\operatorname{Fib}}
\newcommand{\Dep}{\operatorname{Dep}}

\newcommand{\archfixed}{\textsf{schema}}
\newcommand{\instance}{\textsf{instance}}
\newcommand{\param}{\textsf{parameter}}
\newcommand{\state}{\textsf{state}}

\hypersetup{
  pdftitle={Architecture Before the Formula: Receiver-Wise Factorization and the Reconstruction of Transformer Blocks},
  pdfauthor={Luis F. Rosario Freytes},
  pdfsubject={A receiver-factorization reconstruction of Transformer architecture},
  pdfkeywords={Transformers, attention, architecture, representation, receiver-wise factorization, interface, routing, aggregation, local continuation, feedforward network}
}


\title{\textbf{Architecture Before the Formula}\\
\large Receiver-Wise Factorization and the Reconstruction of Transformer Blocks}
\author{Luis F. Rosario Freytes\\University of Michigan, Ann Arbor}
\date{Working preprint v1.0 --- August 2026}

\begin{document}
\maketitle

\begin{abstract}
Neural architectures are usually introduced by formulas, but a represented input--output map does not by itself determine which successor obligations are separately addressed, what information reaches each receiver, or how the resulting local continuation is organized. We define an architecture, conditional on a represented update $F:X\to Y$, as a marked receiver-wise factorization. A marked presentation of the successor state determines receiver branches $B_j$, and an architectural presentation specifies factorizations $B_j=G_jQ_j$, where $Q_j$ is the information made available to receiver $j$ and $G_j$ is its local continuation. This object yields a full-access collapse theorem, a canonical answer quotient that is extensionally sufficient but not architecturally identifying, a refinement preorder on interfaces, and exact and approximate barriers induced by interface fibers. Refining $Q_j$ into retained local state, source-resolved contributions, and aggregation separates what is computed source by source from the receiver macrostate that survives. In an additive linear regime, the refinement produces a state-indexed operator family. Under further commitments---pairwise scalar evidence, positive compositional linking, receiver-row anchoring, source-local payloads, additive aggregation, and one shared finite-separation-rank score family---the receiver interface takes the form of masked query--key softmax attention. In this positive scalar specialization, the routing law also admits receiver-relative effective-potential and free-energy coordinates. The pointwise feedforward sublayer is reconstructed on the other side of the factorization: its marked intermediate activation coordinates refine the receiver-local continuation $G_j$. A canonical PreNorm Transformer block is therefore a coarse attention-constructed receiver interface followed by a structured local continuation, with a finer two-sublayer presentation available when the intermediate residual state is marked. The framework identifies general architectural coordinates, distinguishes provenance typing from realized dependence, and states the additional evidence required to pass from an exact architectural presentation to an implementation-facing mechanism claim. A compact extension allows the marked successor presentation itself to vary across receiver sectors.
\end{abstract}

\paragraph{Assistance disclosure.}
Large language models were used as assistive tools for formalization, literature search, code and typesetting support, and language editing. The author supplied the conceptual program and retains responsibility for verifying the arguments, references, and wording before submission.

\paragraph{Masked-score convention.}
A finite pair score is required only on the admissible receiver--source relation. Hard exclusion is represented by a separate mask. Any finite extension off the admissible relation is mathematically immaterial because masked coefficients are set to zero.

\paragraph{Squared-loss convention.}
Conditional-expectation and squared-loss statements use finite-dimensional Euclidean-valued random targets with finite second moment; every compared predictor is square-integrable on the declared interface.

\section{Architecture as marked receiver-wise factorization}\label{sec:introduction}

A learned formula does not begin from no organization. Its domain and codomain are already represented state spaces; its output may already be resolved into separately addressed obligations; and its computation may already restrict what each obligation can use. Yet the external map alone does not record those distinctions. Two systems can realize the same $F:X\to Y$ while differing in whether one receiver sees the entire predecessor state, a fixed neighborhood, a state-conditioned aggregate, a persistent summary, or a source-resolved family of contributions. Architecture begins where these differences are marked rather than inferred from extensional behavior.

This paper is conditional on a represented update
\[
F:X\longrightarrow Y.
\]
It does not derive represented states, process cuts, or primitive update events. Those belong to a more general theory of represented processes. The narrower question here is:
\begin{quote}
\textbf{What additional marked structure turns a represented update into a receiver-organized architecture, and which restrictions recover the standard Transformer block?}
\end{quote}

The answer is a receiver-wise factorization. A marked presentation of the successor state identifies receiver branches
\[
B_j:X\to W_j.
\]
An architectural presentation then marks, for every receiver $j$, a factorization
\begin{equation}\label{eq:intro-central-factorization}
\boxed{
B_j
=
G_jQ_j,
\qquad
X\xrightarrow{Q_j}Z_j\xrightarrow{G_j}W_j.
}
\end{equation}
The map $Q_j$ is the receiver interface: it determines which predecessor distinctions reach receiver $j$. The map $G_j$ is the receiver-local continuation: it determines what that receiver constructs from the information it receives. Neither map is recoverable uniquely from $B_j$, because arbitrary intermediate factorizations can realize the same branch. They are therefore marked architectural data.

This single factorization organizes the Transformer reconstruction in two directions:
\[
\boxed{
\begin{aligned}
Q_j
&\rightsquigarrow
\text{retained local state, source-resolved routing, attention, and aggregation},\\
G_j
&\rightsquigarrow
\text{residual continuation, normalization, and the pointwise FFN}.
\end{aligned}}
\]
Attention is not taken as a primitive relation or a universal form of state-conditioned computation. It is derived after the transported component of $Q_j$ is restricted to finite source-resolved additive contributions, positive scalar evidence, receiver-row normalization, source-local values, and one shared finite-rank score family. The effective-potential interpretation appears still later, as a logarithmic coordinate on that positive scalar routing law. The FFN is reconstructed differently: its hidden coordinates are architectural because the specified local continuation contains a marked intermediate activation vector, not merely because an arbitrary output map can be written as a formal sum.

\begin{figure}[t]
\centering
\begin{tikzpicture}[
  node distance=12mm and 15mm,
  box/.style={draw,rounded corners,inner sep=5pt,align=center},
  arr/.style={-{Latex[length=2mm]},thick}
]
\node[box] (x) {$X$\\represented predecessor state};
\node[box,right=of x] (q) {$Z_j$\\receiver interface $Q_j$};
\node[box,right=of q] (w) {$W_j$\\receiver obligation};
\draw[arr] (x) -- node[above] {$Q_j$} (q);
\draw[arr] (q) -- node[above] {$G_j$} (w);
\node[box,below=of q] (p) {$\bigl(L_j,\Agg_jP_j\bigr)$\\local state + aggregate exposure};
\draw[arr] (x) |- node[pos=.23,left] {$P_j$} (p);
\draw[arr] (p) -- node[right] {identified with $Z_j$} (q);
\node[above=7mm of q] (b) {$B_j=\pi_j\chi_YF=G_jQ_j$};
\end{tikzpicture}
\caption{The paper's foundational object. Routing and aggregation refine the receiver interface $Q_j$; the FFN refines the local continuation $G_j$.}
\label{fig:central-factorization}
\end{figure}

\subsection{Why the factorization must be marked}

Every branch admits the trivial full-access factorization $B_j=B_j\Id_X$. It also admits a canonical quotient that retains only the equivalence class needed to reproduce that branch. Neither extreme identifies the interface used by the architecture. Likewise, an invertible recoding of the total state may preserve the external map while mixing the coordinates that previously defined receiver-local access. Architecture is therefore not a property of the bare function alone. It is a property of a function together with marked receiver obligations and marked intermediate factorizations.

This is also why the present paper stops before implementation warrant and mechanism identity. An exact factorization may be mathematically valid without descending from an independently specified implementation interface, remaining stable under the scientific operations of interest, or licensing intervention and tracking claims. Establishing those stronger claims requires additional implementation-facing and evidential structure not assumed here.

\subsection{Provenance without a second ontology}

The paper distinguishes the semantic origin of arguments from actual dependence on them. A schema may permit an architectural component
\[
c(B,\xi,\theta,H)
\]
to use schema structure $B$, instance data $\xi$, parameters $\theta$, and current represented state $H$. These argument types are marked. Whether a trained component genuinely varies with one argument over a declared domain is then derived by counterfactual variation. Thus
\[
\text{permitted state dependence}
\neq
\text{realized state dependence}
\neq
\text{variation observed on one restricted window}.
\]
Endogenization replaces a supplied or state-independent slot by a same-typed state-conditioned constructor. Restriction narrows the admissible family of values or maps. The standard Transformer performs both moves: it endogenizes receiver-relative allocation while restricting the construction to a specific pairwise, positive, row-normalized, finite-rank, additive form.

\subsection{Contributions}

The paper makes four connected contributions.

\begin{enumerate}[label=(\arabic*)]
\item \textbf{Receiver-factored architecture.} It defines a complete typed architectural presentation for a represented update, proves that receiver indexing with full shared access collapses to one product-valued map, distinguishes the canonical answer quotient from the marked interface, and orders interfaces by refinement.

\item \textbf{Interface refinement and barriers.} It derives the local-state/exposure split and the source-resolved/aggregate distinction from refinements of $Q_j$. Interface fibers yield exact factorization criteria, deterministic approximation lower bounds, and a squared-loss decomposition into unresolved variation and restricted local accessibility.

\item \textbf{Transformer reconstruction.} It derives additive state-indexed operator families, then recovers masked query--key softmax attention under explicit additional commitments. In the positive scalar regime, it proves the receiver-relative effective-potential and free-energy representation and the log-sum-exp law induced by exact source-carrier quotienting. On the continuation side, it identifies the marked hidden activation coordinates of the FFN and reconstructs the canonical block at coarse and fine architectural grains.

\item \textbf{Architectural coordinates and boundary.} It separates general factorization coordinates from the potential-centered coordinates available only to positive scalar routing, compares neighboring architectures without reducing them to function classes, gives a compact variable-successor extension, and identifies the additional evidence obligations required for stronger implementation-facing and mediation claims.
\end{enumerate}

\subsection{Scope}

The reconstruction is rational rather than historical. It does not claim that Transformers were discovered by successively adding the assumptions listed here, that the standard block is inevitable, or that every state-conditioned interface is attention. It makes no empirical universality claim and does not require experiments to establish its conceptual or mathematical results. Its central claim is conditional: once a represented update and receiver presentation are fixed, the declared factorization identifies the architecture; once further restrictions are imposed on the interface and continuation, the familiar Transformer form follows.

\section{Receiver-factored architectural presentations}\label{sec:receiver-foundation}

This section introduces the minimal object required by the paper. The represented update is assumed. The architecture consists of a marked decomposition of the successor into receiver obligations and a marked factorization of each receiver branch.

\subsection{Represented updates and an upstream barrier}

Fix a represented update
\[
F:X\to Y.
\]
The spaces $X$ and $Y$ are already the states on which the architecture acts. When an upstream representation map
\[
\rho:\Omega\twoheadrightarrow X
\]
is relevant, its fibers impose a prior barrier, but $\rho$ is not part of the architectural primitive studied below.

\begin{lemma}[Fiber factorization criterion]\label{lem:fiber-factorization}
Let $Q:X\twoheadrightarrow Z$ be surjective and let $B:X\to W$. There exists a unique $G:Z\to W$ such that $B=GQ$ if and only if
\[
Q(x)=Q(x')\Longrightarrow B(x)=B(x').
\]
\end{lemma}

\begin{proof}
If $B=GQ$, equality of interface values implies equality of branch values. Conversely, if $B$ is constant on every $Q$-fiber, define $G(z)=B(x)$ for any $x$ with $Q(x)=z$. The fiber condition makes $G$ well defined, and surjectivity gives uniqueness.
\end{proof}

\begin{corollary}[Upstream representation barrier]\label{cor:representation-barrier}
A target $f:\Omega\to W$ factors through $\rho$ if and only if it is constant on the fibers of $\rho$. No downstream architecture operating only on $X$ can recover distinctions erased by $\rho$.
\end{corollary}

The remainder of the paper is conditional on $X$ and begins after this upstream identification has already occurred.

\subsection{Marked receiver presentations}

\begin{definition}[Marked receiver presentation]\label{def:marked-receiver-presentation}
A marked receiver presentation of $Y$ consists of:
\begin{enumerate}[label=(\roman*)]
\item a finite receiver set $R$;
\item receiver value spaces $(W_j)_{j\in R}$;
\item a bijection
\[
\chi_Y:Y\xrightarrow{\cong}H_Y\subseteq\prod_{j\in R}W_j;
\]
\item distinguished receiver projections $\pi_j:H_Y\to W_j$.
\end{enumerate}
\end{definition}

The bijection makes this a complete presentation of the represented successor state rather than a further quotient of it. If a scientific question intentionally replaces $Y$ by a quotient $\bar Y$, the definition is applied to the induced update $\bar F:X\to\bar Y$. A noninjective map out of the original $Y$ would instead specify only a partial or quotient target profile and would not be a complete receiver presentation of $F$.

The presentation is marked because the projections are part of the architectural claim. A global recoding of $Y$ can preserve total-state information while destroying the identification of one coordinate as receiver $j$.

\begin{definition}[Receiver branch]\label{def:receiver-branch}
For receiver $j\in R$, define
\[
B_j:=\pi_j\chi_YF:X\to W_j.
\]
\end{definition}

\begin{proposition}[Full-access receiver collapse]\label{prop:full-access-collapse}
Every branch has the trivial factorization $B_j=B_j\Id_X$. Moreover, the family $(B_j)_{j\in R}$ is uniquely equivalent to the product-valued map
\[
B:=\langle B_j\rangle_{j\in R}:X\to\prod_{j\in R}W_j,
\]
whose image lies in $H_Y$. Because $\chi_Y$ is bijective, $F=\chi_Y^{-1}B$. Thus receiver labels with complete shared access impose no nontrivial access architecture beyond the marked successor presentation itself.
\end{proposition}

\begin{proof}
The first claim is immediate. The product map is defined by $B(x)=(B_j(x))_j$ and is uniquely determined by its projections. By construction, $B=\chi_YF$, so its image lies in $H_Y$ and applying $\chi_Y^{-1}$ recovers $F$.
\end{proof}

\subsection{The architectural object}

\begin{definition}[Receiver-factored architectural presentation]\label{def:receiver-factored-architecture}
Let $F:X\to Y$ have the marked receiver presentation of \Cref{def:marked-receiver-presentation}. A receiver-factored architectural presentation is
\[
\mathcal A
=
\left(
F,\chi_Y,
\{X\xrightarrow{Q_j}Z_j\xrightarrow{G_j}W_j\}_{j\in R}
\right),
\]
where each $Q_j:X\twoheadrightarrow Z_j$ is surjective and
\begin{equation}\label{eq:receiver-factorization}
\boxed{
B_j
=
\pi_j\chi_YF
=
G_jQ_j.
}
\end{equation}
The map $Q_j$ is the receiver interface and $G_j$ is the receiver-local continuation.
\end{definition}

Every field in this object is typed and used in \eqref{eq:receiver-factorization}. The definition does not add a catch-all space of messages, graphs, procedures, or presentations. Those appear only when a particular $Q_j$ or $G_j$ is refined.

\begin{proposition}[Receiver-interface re-presentation]\label{prop:interface-representation}
Let $B_j=G_jQ_j$ with $Q_j:X\twoheadrightarrow Z_j$, and let $\zeta_j:Z_j\to Z'_j$ be a bijection. Define
\[
Q'_j:=\zeta_jQ_j,
\qquad
G'_j:=G_j\zeta_j^{-1}.
\]
Then $B_j=G'_jQ'_j$, and $Q_j$ and $Q'_j$ induce the same partition of $X$. We treat them as coordinate re-presentations of the same receiver interface.
\end{proposition}

\begin{proof}
The branch identity follows by cancellation of $\zeta_j^{-1}\zeta_j$. Because $\zeta_j$ is injective,
\[
Q'_j(x)=Q'_j(x')
\quad\Longleftrightarrow\quad
Q_j(x)=Q_j(x'),
\]
so the two interfaces have exactly the same fibers.
\end{proof}

This convention concerns only bijective recoding of one receiver interface. Comparisons that also transport receiver systems, schedules, interventions, implementation witnesses, or tracked roles require additional marked structure beyond the present paper.

\begin{definition}[Receiver indistinguishability]\label{def:receiver-indistinguishability}
For receiver $j$, define
\[
x\sim_{Q_j}x'
\quad\Longleftrightarrow\quad
Q_j(x)=Q_j(x').
\]
\end{definition}

By \Cref{lem:fiber-factorization}, every continuation through $Q_j$ is constant on these equivalence classes. Receiver locality is therefore exact relative to the marked interface: it is a declaration of which predecessor distinctions can affect that branch through the specified continuation.

\subsection{The canonical answer quotient is not the architecture}

For one branch $B:X\to W$, define
\[
x\sim_Bx'
\quad\Longleftrightarrow\quad
B(x)=B(x')
\]
and let $q_B:X\twoheadrightarrow X/{\sim_B}$ be the quotient.

\begin{proposition}[Canonical exact branch quotient]\label{prop:canonical-branch-quotient}
The branch $B$ factors uniquely through $q_B$. If $Q:X\twoheadrightarrow Z$ is any exact receiver interface for $B$, then there is a unique $h:Z\to X/{\sim_B}$ such that
\[
q_B=hQ.
\]
Hence $q_B$ is the coarsest exact quotient sufficient for the selected branch.
\end{proposition}

\begin{proof}
Define $\bar B([x])=B(x)$. If $B=GQ$, define $h(Q(x))=[x]$. Exactness implies that $Q(x)=Q(x')$ only when $B(x)=B(x')$, so $h$ is well defined. Surjectivity gives uniqueness.
\end{proof}

The two extensional extremes are
\[
Q=\Id_X
\qquad\text{and}\qquad
Q=q_B.
\]
The first retains everything; the second retains only the answer-equivalence class. Neither identifies the intermediate organization actually used. A receiver interface is therefore marked independently of the branch it realizes.

\subsection{Interface refinement}

\begin{definition}[Interface refinement]\label{def:interface-refinement}
For surjective interfaces $Q:X\twoheadrightarrow Z$ and $Q':X\twoheadrightarrow Z'$, write
\[
Q\preceq Q'
\]
when there is a map $h:Z'\to Z$ with $Q=hQ'$. Then $Q'$ retains every distinction retained by $Q$ and possibly more.
\end{definition}

\begin{proposition}[Refinement preorder and branch monotonicity]\label{prop:refinement-preorder}
The relation $\preceq$ is reflexive and transitive. If $Q\preceq Q'$ and a branch factors through $Q$, then it also factors through $Q'$.
\end{proposition}

\begin{proof}
Reflexivity uses the identity; transitivity uses composition. If $Q=hQ'$ and $B=GQ$, then $B=(Gh)Q'$.
\end{proof}

\subsection{Locality is presentation-relative}

\begin{proposition}[Global conjugacy can destroy coordinate locality]\label{prop:locality-conjugacy}
Let $Y=\R^2$ with marked receiver coordinates and define
\[
F(x,y)=(x^2,y),
\qquad
T(x,y)=(x+y,x-y).
\]
The first branch of $F$ depends only on $x$, but both branches of $TFT^{-1}$ depend on both mixed coordinates.
\end{proposition}

\begin{proof}
Since
\[
T^{-1}(u,v)=\left(\frac{u+v}{2},\frac{u-v}{2}\right),
\]
one obtains
\[
TFT^{-1}(u,v)
=
\left(
\frac{(u+v)^2}{4}+\frac{u-v}{2},
\frac{(u+v)^2}{4}-\frac{u-v}{2}
\right).
\]
Both components depend on $u$ and $v$.
\end{proof}

Functional equivalence under an invertible total-state recoding does not preserve the receiver marks or interfaces. Architectural locality belongs to the marked presentation, not to the unmarked function class.

\section{Refining receiver interfaces: source resolution, aggregation, and barriers}\label{sec:interface-refinement}

The architectural object of \Cref{def:receiver-factored-architecture} does not assume that every interface has the same internal form. This section introduces two successive refinements that are central to residual routed architectures:
\[
Q_j
\rightsquigarrow
\langle L_j,E_j\rangle
\rightsquigarrow
\left\langle L_j,\Agg_jP_j\right\rangle.
\]
The first distinguishes receiver-local state from transported exposure. The second distinguishes source-resolved computation from the macrostate retained after aggregation.

\subsection{Retained local state and transported exposure}

\begin{definition}[Local-state/exposure refinement]\label{def:local-exposure-refinement}
A receiver interface $Q_j:X\twoheadrightarrow Z_j$ has a local-state/exposure refinement when there are maps
\[
L_j:X\to U_j,
\qquad
E_j:X\to Z_j^{\mathrm{exp}},
\]
and a bijection from the realized image of $\langle L_j,E_j\rangle$ to $Z_j$ under which $Q_j$ is identified with
\[
Q_j\cong\langle L_j,E_j\rangle.
\]
If no separately retained local state is marked, take $U_j$ to be a singleton.
\end{definition}

For a same-carrier residual block, $L_i(H)=h_i$. The exposure $E_i(H)$ is the additional state transported or constructed from the represented source family. The distinction is architectural: the same aggregate may lead to a different successor when the bypassed local state differs.

\subsection{Source-resolved refinement}

\begin{definition}[Source-resolved receiver refinement]\label{def:source-resolved-refinement}
The transported exposure $E_j:X\to Z_j^{\mathrm{exp}}$ has a source-resolved refinement when there are
\[
P_j:X\to\mathcal M_j,
\qquad
\Agg_j:\mathcal M_j\to Z_j^{\mathrm{exp}}
\]
such that
\[
E_j=\Agg_jP_j.
\]
For a finite source set $S$, a common realized presentation is
\[
P_j(H)=\bigl(m_{ji}(H)\bigr)_{i\in S},
\]
where each $m_{ji}(H)$ is a typed contribution to receiver $j$ from source $i$.
\end{definition}

The three objects
\[
P_j(H),
\qquad
\Agg_jP_j(H),
\qquad
Q_j(H)\cong\bigl(L_j(H),\Agg_jP_j(H)\bigr)
\]
must remain distinct. The first preserves source resolution. The second is the transported receiver macrostate. The third is the complete information available to the receiver-local continuation.

\begin{figure}[t]
\centering
\begin{tikzpicture}[
  node distance=10mm and 12mm,
  box/.style={draw,rounded corners,inner sep=4pt,align=center},
  arr/.style={-{Latex[length=2mm]},thick}
]
\node[box] (x) {$X$};
\node[box,right=of x] (p) {$\mathcal M_j$\\source-resolved $P_j$};
\node[box,right=of p] (e) {$Z_j^{\mathrm{exp}}$\\aggregate exposure};
\node[box,right=of e] (q) {$Z_j$\\complete interface};
\node[box,right=of q] (w) {$W_j$};
\draw[arr] (x) -- node[above] {$P_j$} (p);
\draw[arr] (p) -- node[above] {$\Agg_j$} (e);
\draw[arr] (e) -- node[above] {combine with $L_j$} (q);
\draw[arr] (q) -- node[above] {$G_j$} (w);
\draw[arr] (x) to[bend right=34] node[below] {$L_j$} (q);
\end{tikzpicture}
\caption{Source resolution, aggregation, and the complete receiver interface. Every many-to-one arrow creates a possible identification barrier for later continuations.}
\label{fig:interface-refinement}
\end{figure}

\subsection{A concrete aggregation collision}

\begin{example}[A moment can erase source-resolved variation]\label{ex:aggregation-collision}
Let
\[
P(x_1,x_2)=(x_1,x_2),
\qquad
\Agg(x_1,x_2)=x_1+x_2.
\]
Then $(1,-1)$ and $(0,0)$ have the same aggregate, but the source-resolved target
\[
f(x_1,x_2)=x_1^2+x_2^2
\]
takes values $2$ and $0$. No continuation that receives only $x_1+x_2$ can reproduce $f$ on both inputs.
\end{example}

The example is not a criticism of aggregation. A sum may be the intended sufficient statistic. It identifies the exact question: which target distinctions vary inside the fibers of the retained exposure?

\subsection{Exact and approximate barriers at the aggregate cut}

Let
\[
P:X\to\mathcal M,
\qquad
\Agg:\mathcal M\to Z,
\qquad
E:=\Agg P.
\]
Corestrict $E$ to its realized image so that $E:X\twoheadrightarrow Z$ is surjective.

\begin{theorem}[Exact aggregation-fiber obstruction]\label{thm:aggregation-fiber}
For every continuation $G:Z\to W$,
\[
E(x)=E(x')
\Longrightarrow
G(E(x))=G(E(x')).
\]
No continuation after the aggregate cut can distinguish predecessor states identified by $E$.
\end{theorem}

\begin{proof}
Apply the same function $G$ to the equal interface values.
\end{proof}

The theorem is elementary but architecturally decisive. A receiver may compute a large source-resolved family while retaining only one fixed-dimensional statistic. Computation breadth does not imply source-resolved access by the next local continuation.

\begin{definition}[Target variation along interface fibers]\label{def:fiber-variation}
Let $(W,d)$ be a metric space and $f:X\to W$. Define
\[
\Delta_E(f)
:=
\sup_{z\in Z}\diam f(E^{-1}(z))
=
\sup_{E(x)=E(x')}d(f(x),f(x')).
\]
\end{definition}

\begin{theorem}[Deterministic approximation lower bound]\label{thm:deterministic-lower-bound}
For every $G:Z\to W$,
\[
\boxed{
\sup_{x\in X}d\bigl(f(x),G(E(x))\bigr)
\ge
\frac12\Delta_E(f).
}
\]
\end{theorem}

\begin{proof}
Fix $x,x'$ with $E(x)=E(x')=:z$. By the triangle inequality,
\[
d(f(x),f(x'))
\le d(f(x),G(z))+d(G(z),f(x')).
\]
At least one term on the right is at least half the left side. Take the supremum over common-fiber pairs.
\end{proof}

\begin{corollary}[Exact recoverability criterion]\label{cor:exact-recoverability}
There exists $G:Z\to W$ with $f=GE$ if and only if $\Delta_E(f)=0$.
\end{corollary}

\begin{proof}
The forward direction is immediate. The reverse direction follows from \Cref{lem:fiber-factorization} because $f$ is constant on every $E$-fiber.
\end{proof}

\begin{proposition}[Aggregation cannot improve target resolution]\label{prop:aggregation-monotonicity}
If $E=\Agg P$, then
\[
\Delta_E(f)\ge\Delta_P(f).
\]
\end{proposition}

\begin{proof}
Every $P$-fiber is contained in an $E$-fiber, so taking the supremum over the coarser $E$-fibers cannot decrease the target diameter.
\end{proof}

These statements separate interface revision from continuation revision. Enlarging $G$ can improve how the retained interface is used; it cannot recover a distinction absent from $E$.

\subsection{Unresolved variation and restricted local accessibility}

Exact fibers isolate information availability. A separate defect arises when the chosen local continuation family cannot use all distinctions retained by the interface.

Let $(H,Y)$ be square-integrable random variables, let $Z=E(H)$, and let $\mathcal G$ be a class of measurable maps $g$ with $\E\norm{g(Z)}^2<\infty$.

\begin{theorem}[Squared-loss unresolved-variation/accessibility decomposition]\label{thm:squared-decomposition}
Let $m(Z):=\E[Y\mid Z]$. Then
\[
\boxed{
\inf_{g\in\mathcal G}\E\norm{Y-g(Z)}^2
=
\E\norm{Y-m(Z)}^2
+
\inf_{g\in\mathcal G}\E\norm{m(Z)-g(Z)}^2.
}
\]
\end{theorem}

\begin{proof}
For any $g\in\mathcal G$,
\[
Y-g(Z)=(Y-m(Z))+(m(Z)-g(Z)).
\]
The cross term in the squared norm has expectation zero because $m(Z)-g(Z)$ is $Z$-measurable and $\E[Y-m(Z)\mid Z]=0$. Taking the infimum proves the identity.
\end{proof}

The first term is irreducible given the receiver interface: target variation remains unresolved inside its fibers. The second term is an accessibility defect of the declared local family. In a Transformer, enlarging the FFN may reduce the second term but not the first. This is the precise sense in which the interface and the local continuation are different architectural resources.

\subsection{Information accounting}

For finite-valued random variables, a deterministic quotient also admits the exact identity
\[
I(Y;P(H))
=
I(Y;E(H))
+
I\bigl(Y;P(H)\mid E(H)\bigr).
\]
The conditional term measures target-relevant evidence present in the source-resolved presentation but absent from the retained aggregate. The proof is the mutual-information chain rule and is recorded in the technical supplement. The distribution-free fiber theorems remain primary because they do not depend on a probability law or a selected target distribution.

\section{Provenance, endogenization, and routed realizations}\label{sec:provenance-routing}

The receiver interface specifies what reaches a continuation. A second question asks where the maps that construct that interface obtain their operative values. The answer requires semantic argument types, but it does not require a second ontology of state-indexed presentations. The actual routed law is given by the evaluated maps that refine $Q_j$.

\subsection{Marked argument origins and derived dependence}

Let an architectural component have a typed signature
\begin{equation}\label{eq:typed-component-signature}
c:
D_B\times D_\xi\times D_\theta\times X
\longrightarrow C.
\end{equation}
The arguments have declared meanings:
\[
\begin{aligned}
B&:\text{schema-declared structure},
&\xi&:\text{instance-supplied structure},\\
\theta&:\text{parameter data},
&H\in X&:\text{current represented state}.
\end{aligned}
\]
These semantic source types are marked. Actual dependence is derived relative to an admissible domain.

\begin{definition}[Essential dependence]\label{def:essential-dependence}
Let
\[
D\subseteq D_B\times D_\xi\times D_\theta\times X
\]
be an admissible domain. For $a\in\{B,\xi,\theta,H\}$, write
\[
a\in\Dep_D(c)
\]
when there exist $u,v\in D$ that agree in every coordinate except possibly $a$ and satisfy $c(u)\neq c(v)$.
\end{definition}

The definition distinguishes three claims:
\[
\boxed{
\begin{aligned}
\text{the schema permits an }H\text{-argument}
&\neq H\in\Dep_D(c),\\
H\in\Dep_D(c)
&\neq H\in\Dep_K(c)
\quad\text{for every observed }K\subseteq D.
\end{aligned}}
\]
A trained constructor can be state-dependent on its admissible domain while remaining constant on the states observed in one dataset or deployment window.

\begin{proposition}[Observed behavior underdetermines state provenance]\label{prop:provenance-underdetermination}
Let $\varnothing\neq K\subsetneq X$ and let $A_0\neq A_1$. Define
\[
\Phi_0(H)=A_0,
\qquad
\Phi_1(H)=
\begin{cases}
A_0,&H\in K,\\
A_1,&H\notin K.
\end{cases}
\]
Then $H\notin\Dep_X(\Phi_0)$ and $H\in\Dep_X(\Phi_1)$, while $\Phi_0=\Phi_1$ on $K$.
\end{proposition}

\begin{proof}
The first map is constant. The second takes two values because both $K$ and its complement are nonempty. Their agreement on $K$ is immediate.
\end{proof}

The extensional routing law on a restricted support therefore cannot establish whether its operative relation was supplied, learned but state-independent, or generated from state and accidentally constant there. The argument origins must be marked, and realized dependence must be tested on a declared admissible domain.

\subsection{Restriction and endogenization}

\begin{definition}[Restriction]\label{def:restriction-move}
A restriction replaces an admissible family $\Sigma$ for one typed architectural component by a subfamily
\[
\Sigma'\subseteq\Sigma.
\]
\end{definition}

\begin{definition}[Endogenization]\label{def:endogenization-move}
An endogenization replaces a state-independent component
\[
c(B,\xi,\theta)
\]
by a same-typed constructor
\[
\widetilde c(B,\xi,\theta,H).
\]
\end{definition}

Restriction changes which values or maps are admissible. Endogenization changes which arguments may construct the operative value. They are orthogonal. A fixed graph can carry state-dependent weights; a state-generated graph can still be restricted to bounded degree; a state-conditioned coefficient can remain constant on the realized domain.

The standard Transformer illustrates both moves. Relative to a supplied relation, it endogenizes receiver-relative allocation by making it a function of the current token states. It simultaneously restricts the construction to pairwise scalar evidence, positive support, receiver-row normalization, source-local values, additive aggregation, and a shared finite-rank score family.

\subsection{Operator-valued pair contributions}

Assume now that the source-resolved exposure of receiver $r$ is additive and vector-valued. Let $R$ be a finite receiver set and $S$ a finite source set. For each $i\in S$, let $U_i$ be a real vector space, and for each $r\in R$, let $W_r$ be a real receiver-exposure vector space.

\begin{definition}[Additive routed realization]\label{def:additive-routed-realization}
An additive routed realization of $E_r$ consists of payload maps
\[
u_i:X\to U_i
\]
and pair-contribution operators
\[
C_{ri}:X\to\mathcal L(U_i,W_r)
\]
such that
\begin{equation}\label{eq:additive-route}
\boxed{
E_r(H)
=
\sum_{i\in S}C_{ri}(H)u_i(H).
}
\end{equation}
The corresponding source-resolved presentation is
\[
P_r(H)=\bigl(C_{ri}(H)u_i(H)\bigr)_{i\in S},
\]
and $\Agg_r$ is summation.
\end{definition}

The operator $C_{ri}(H)$ is the minimal pair object in the additive linear regime. A further factorization
\begin{equation}\label{eq:scalar-transport-factorization}
C_{ri}(H)=\lambda_{ri}(H)T_{ri}(H)
\end{equation}
becomes architecturally meaningful only when scalar relevance and payload transport are separately marked components. Without that mark, the rescaling
\[
\lambda_{ri}\mapsto a_{ri}\lambda_{ri},
\qquad
T_{ri}\mapsto a_{ri}^{-1}T_{ri}
\]
leaves the pair operator unchanged wherever $a_{ri}\neq0$.

For fixed $H$, define
\begin{equation}\label{eq:state-operator}
A_H:\bigoplus_{i\in S}U_i\to\bigoplus_{r\in R}W_r,
\qquad
(A_Hu)_r:=\sum_iC_{ri}(H)u_i.
\end{equation}

\begin{proposition}[Routed interfaces define state-indexed operator families]\label{prop:routed-operator-field}
For every fixed $H$, $A_H$ is linear and
\[
E(H)=A_Hu(H).
\]
Thus the routed realization determines a state-to-operator map $H\mapsto A_H$.
\end{proposition}

\begin{proof}
Linearity follows from the linearity of every $C_{ri}(H)$ and of finite summation. Substituting the state-derived payload family gives \eqref{eq:additive-route} receiver by receiver.
\end{proof}

The current state performs two distinct roles: it supplies the payload $u(H)$ and may also construct the operator $A_H$ that acts on that payload. This state-to-operator lift is broader than attention.

\subsection{Support, contribution, and retained exposure}

Define the operative pair support
\[
\mathsf{Supp}_H:=\{(r,i):C_{ri}(H)\neq0\}.
\]
The following distinctions are structural:
\[
\boxed{
\begin{gathered}
\text{source carrier}
\neq
\text{admissible pair relation}
\neq
\text{operative support},\\
\text{pair operator}
\neq
\text{source payload}
\neq
\text{retained aggregate exposure}.
\end{gathered}}
\]
A nonzero pair operator may act on a zero payload. Several nonzero source contributions may cancel under summation. Different pair families may induce the same aggregate. These are reasons not to identify a routing matrix with direct mediation or with the complete information available to a receiver.

\subsection{Controlled contrasts}

The factorization coordinates place several familiar architectures without treating them as historical steps toward the Transformer.

\paragraph{Linear prediction.}
A linear predictor $x\mapsto Wx+b$ acts after a representation has already fixed $x$. At the one-receiver grain, the interface may be all of $x$ and the continuation is affine. The model does not by itself construct a receiver-relative source organization.

\paragraph{Normalized kernel estimation.}
A normalized kernel estimator
\[
\widehat y(x)
=
\sum_i
\frac{\kappa(x,x_i)}{\sum_j\kappa(x,x_j)}y_i
\]
already has a one-receiver source-resolved additive interface. Its allocation is query-conditioned, but the comparison geometry is normally supplied or globally parameterized and the result is usually one terminal observable rather than a successor representation that constructs a new routed law \citep{nadaraya1964estimating,watson1964smooth}.

\paragraph{Convolution.}
For a finite offset family $D$,
\[
(\operatorname{Conv}H)_r
=
\sum_{\delta\in D}W_\delta h_{r+\delta}.
\]
Here $C_{r,r+\delta}=W_\delta$ on a schema-fixed translation-generated support and zero elsewhere. The relation and parameter sharing are supplied by the schema; the payloads vary with the state.

\paragraph{Graph message passing.}
A graph layer
\[
E_r(H,\xi)
=
\Agg_{i\in N_\xi(r)}m_\theta(h_r,h_i,e_{ri})
\]
uses an instance-supplied admissible relation while allowing state-dependent messages or weights. Graph attention endogenizes allocation inside the supplied accessible sector rather than eliminating the graph \citep{velickovic2018graph,bronstein2021geometric}.

\paragraph{Recurrence and state-space summaries.}
A recurrent architecture gives the receiver a persistent summary instead of a source-resolved row over all predecessors. Long-range influence is mediated through that summary, so the interface has different fibers even when the external function can be matched.

\paragraph{Mixture-of-experts routing.}
A token state can construct coefficients over a supplied expert carrier. This is a routed interface, but its source and receiver types differ from token self-attention and its capacity constraints may couple allocations across tokens.

No scalar measure globally orders these systems. One may endogenize support while fixing aggregation, enlarge message types while compressing source resolution, or vary carrier populations while retaining a rigid schedule. The next section selects a narrower additive regime and derives the additional commitments that make it attention.

\section{Attention as a restricted receiver-interface realization}\label{sec:attention}

The state-indexed operator family of \Cref{sec:provenance-routing} does not imply attention. Attention is obtained by restricting the source-resolved component of $Q_r$ to a particular scalar, positive, row-anchored, source-local, additive, and finitely separable form. The effective-potential representation appears only after that routing row has been derived.

The assumptions below select this narrower regime rather than axiomatizing receiver-organized computation in general. Their distinct architectural work is exhibited by explicit neighboring constructions in the technical supplement; removing one commitment changes a typed component of the realization rather than making the surrounding factorization meaningless.

Let $R$ be a finite receiver set, let $S$ be a finite source set, and let $H\in X$ be the current represented state. Fix typed receiver and source feature maps
\[
x_r:X\to\mathcal X_R,
\qquad
y_i:X\to\mathcal X_S.
\]
For self-attention, $R=S$ and one usually has $x_i(H)=y_i(H)=h_i$. Keeping the two feature families distinct makes the normal form equally well typed for cross-attention and other receiver--source architectures.

\subsection{Commitments before the formula}

\begin{assumption}[Pairwise scalar evidence]\label{ass:pairwise-scalar}
For every admissible pair $(r,i)$, the routed interface uses a finite scalar evidence coordinate
\[
s_{ri}(H)\in\R.
\]
Hard exclusion is represented separately by a mask.
\end{assumption}

\begin{assumption}[Positive baseline and nonempty rows]\label{ass:positive-baseline}
There is a hard mask $M_{ri}\in\{0,1\}$ and a positive baseline factor $\kappa_{ri}>0$ for every pair. For each receiver,
\[
S_r:=\{i\in S:M_{ri}=1\}
\]
is nonempty. Values of $\kappa_{ri}$ off the admissible support are immaterial.
\end{assumption}

The single factor $\kappa_{ri}$ contains every state-independent positive pair term that is extensionally relevant to the row. It may be decomposed into schema-, instance-, or parameter-origin factors only when those factors are independently marked. A finite additive pair bias is equivalent to multiplying $\kappa_{ri}$ by its exponential and is not counted twice.

\begin{assumption}[Compositional positive link]\label{ass:link}
There is a continuous strictly increasing map
\[
\psi:\R\to(0,\infty)
\]
with
\[
\psi(a+b)=\psi(a)\psi(b).
\]
Thus independent additive evidence coordinates combine multiplicatively as positive raw support.
\end{assumption}

\begin{assumption}[Receiver-row anchoring]\label{ass:row-anchor}
For each receiver $r$, positive raw weights $(w_{ri})_{i\in S_r}$ are converted to coefficients $(\alpha_{ri})_{i\in S}$ such that:
\begin{enumerate}[label=(\roman*)]
\item $\alpha_{ri}=0$ for $i\notin S_r$;
\item $\alpha_{ri}>0$ for $i\in S_r$ and $\sum_{i\in S_r}\alpha_{ri}=1$;
\item within-row ratios are preserved:
\[
\frac{\alpha_{ri}}{\alpha_{rj}}
=
\frac{w_{ri}}{w_{rj}}
\qquad(i,j\in S_r).
\]
\end{enumerate}
\end{assumption}

\begin{assumption}[Source-local payload and additive aggregate]\label{ass:values-aggregation}
Each source has a payload map
\[
v_i:\mathcal X_S\to W,
\]
so its evaluated payload $v_i(y_i(H))$ factors through the marked source feature and is independent of the receiver. The transported exposure is
\[
E_r(H)=z_r(H):=\sum_{i\in S}\alpha_{ri}(H)v_i(y_i(H)).
\]
\end{assumption}

The commitment chain is
\[
\boxed{
Q_r
\rightsquigarrow
P_r
\rightsquigarrow
\bigl(s_{ri},\kappa_{ri},M_{ri}\bigr)_i
\rightsquigarrow
\alpha_r
\rightsquigarrow
z_r.
}
\]
More general receiver interfaces may use higher-order contributions, signed or operator-valued coefficients, pair-conditioned payloads, unnormalized flow, globally coupled plans, recurrent aggregation, or retained source-resolved memory.

\subsection{The positive link and row anchor}

\begin{lemma}[Exponential evidence link]\label{lem:exponential-link}
Under \Cref{ass:link}, there exists $\beta>0$ such that
\[
\psi(t)=e^{\beta t}.
\]
\end{lemma}

\begin{proof}
The multiplicative law gives $\psi(0)=1$. The continuous map $\varphi=\log\psi$ satisfies the additive Cauchy equation, so $\varphi(t)=\beta t$. Strict monotonicity implies $\beta>0$.
\end{proof}

Set $\tau:=\beta^{-1}>0$. This is a score scale controlling the evidence--dispersion tradeoff; no physical temperature claim follows.

\begin{lemma}[Uniqueness of receiver-row anchoring]\label{lem:row-anchor}
Let $(w_i)_{i\in A}$ be positive on a nonempty finite set. The unique positive unit-mass vector preserving all ratios $w_i/w_j$ is
\[
\alpha_i=\frac{w_i}{\sum_{j\in A}w_j}.
\]
\end{lemma}

\begin{proof}
Ratio preservation implies $\alpha_i=cw_i$ for one $c>0$. Unit mass determines $c=(\sum_jw_j)^{-1}$.
\end{proof}

Thus, on admissible pairs,
\[
w_{ri}(H)=\kappa_{ri}e^{s_{ri}(H)/\tau}
\]
and
\begin{equation}\label{eq:general-softmax-row}
\alpha_{ri}(H)
=
\frac{M_{ri}\kappa_{ri}e^{s_{ri}(H)/\tau}}
{\sum_{j\in S_r}\kappa_{rj}e^{s_{rj}(H)/\tau}}.
\end{equation}
Row normalization is therefore not a primitive truth of architecture. It follows from positivity, receiver-local unit mass, and preservation of relative evidence \citep{luce1959individual}.

\subsection{Shared finite separation rank}

Low rank of one realized score matrix is not enough to produce shared query and key maps. The commitment concerns the score family across admissible states.

\begin{definition}[Shared finite separation rank]\label{def:shared-separation-rank}
The contextual score family has shared separation rank at most $d$ when there exist functions
\[
f_1,\ldots,f_d:\mathcal X_R\to\R,
\qquad
g_1,\ldots,g_d:\mathcal X_S\to\R
\]
such that, for every admissible pair $(r,i)$ and represented state $H$,
\begin{equation}\label{eq:separable-score}
s_{ri}(H)
=
\sum_{\ell=1}^{d}f_\ell(x_r(H))g_\ell(y_i(H)).
\end{equation}
The same feature functions are used across receivers, sources, and states.
\end{definition}

\begin{proposition}[Query--key coordinates from finite separation rank]\label{prop:separable-bilinear}
A contextual score family has the form \eqref{eq:separable-score} if and only if there are maps
\[
q:\mathcal X_R\to\R^d,
\qquad
k:\mathcal X_S\to\R^d
\]
with
\[
s_{ri}(H)=\ip{q(x_r(H))}{k(y_i(H))}
\]
for every admissible pair and represented state.
\end{proposition}

\begin{proof}
Given \eqref{eq:separable-score}, set $q(x)=(f_\ell(x))_\ell$ and $k(y)=(g_\ell(y))_\ell$. Evaluating these maps at $x_r(H)$ and $y_i(H)$ gives the stated score. The converse follows by expanding the Euclidean inner product coordinatewise.
\end{proof}

The query--key coordinates are therefore derived coordinates on one shared separable score family. They are not semantic primitives and are not uniquely determined.

\subsection{Attention normal form}

\begin{theorem}[Single-head attention normal form]\label{thm:attention-normal-form}
Assume \Cref{ass:pairwise-scalar,ass:positive-baseline,ass:link,ass:row-anchor,ass:values-aggregation} and shared separation rank at most $d$. Then there are shared query and key maps $q,k$ such that
\begin{equation}\label{eq:attention-normal-form}
\boxed{
z_r(H)=\sum_{i\in S}\alpha_{ri}(H)v_i(y_i(H)),
}
\end{equation}
where
\begin{equation}\label{eq:attention-weights}
\boxed{
\alpha_{ri}(H)
=
\frac{
M_{ri}\kappa_{ri}
\exp\!\left(\ip{q(x_r(H))}{k(y_i(H))}/\tau\right)
}{
\sum_{j\in S_r}\kappa_{rj}
\exp\!\left(\ip{q(x_r(H))}{k(y_j(H))}/\tau\right)
}.
}
\end{equation}
\end{theorem}

\begin{proof}
The exponential-link lemma gives positive raw weights $\kappa_{ri}e^{s_{ri}(H)/\tau}$. The separation-rank proposition gives $s_{ri}(H)=\ip{q(x_r(H))}{k(y_i(H))}$. The row-anchor lemma gives \eqref{eq:attention-weights}, and source-local payloads with additive aggregation give \eqref{eq:attention-normal-form}.
\end{proof}

\begin{corollary}[Standard projected attention]\label{cor:standard-attention}
In self-attention, if $R=S$, $x_i(H)=y_i(H)=h_i$, and
\[
q(h)=W_Qh,
\qquad
k(h)=W_Kh,
\qquad
v_i(h)=W_Vh,
\]
and a finite supplied pair bias $b_{ri}$ is represented by $\kappa_{ri}=e^{b_{ri}/\tau}$ on admissible pairs, then \Cref{thm:attention-normal-form} is masked scaled dot-product attention after choosing the conventional score scale \citep{vaswani2017attention}.
\end{corollary}

A state-dependent pair bias should instead remain in $s_{ri}(H)$. The decomposition between $\kappa$ and $s$ is provenance-relative: it is meaningful when the architecture marks which term is state-independent and which is constructed from current state.

\begin{proposition}[Converse verification]\label{prop:attention-converse}
A standard masked single-head attention layer with at least one admissible source per receiver satisfies the assumptions above. Hence its receiver exposure has both the routed normal form and the effective-potential representation derived below.
\end{proposition}

\begin{proof}
The layer uses finite bilinear scores on admissible pairs, positive exponentiation, rowwise normalization, source-local values, and additive aggregation. A finite mask bias can be represented by $\kappa$; forbidden pairs are represented by $M=0$.
\end{proof}

The theorem and converse classify a declared regime. They do not imply that every state-conditioned interface, every graph constructor, or every operator field is attention.

\subsection{Effective potential after the routing law}\label{subsec:effective-potential}

Normalize the baseline on each admissible row:
\begin{equation}\label{eq:reference-measure}
\pi_{ri}:=\frac{\kappa_{ri}}{\sum_{j\in S_r}\kappa_{rj}}
\qquad(i\in S_r).
\end{equation}
Then $\pi_r$ is a reference probability measure on $S_r$.

\begin{definition}[Receiver-relative effective potential]\label{def:effective-potential}
For $i\in S_r$, define
\begin{equation}\label{eq:effective-potential}
\mathcal U_{ri}(H):=-s_{ri}(H)-\tau\log\pi_{ri}.
\end{equation}
For $i\notin S_r$, set $\mathcal U_{ri}(H)=+\infty$.
\end{definition}

Then
\begin{equation}\label{eq:gibbs-row}
\alpha_{ri}(H)
=
\frac{e^{-\mathcal U_{ri}(H)/\tau}}
{\sum_{j\in S_r}e^{-\mathcal U_{rj}(H)/\tau}}.
\end{equation}
The mask chooses the accessible source sector, the baseline supplies the state-independent reference landscape, and current-state evidence deforms that landscape. This is an exact coordinate on the positive scalar routing law, not the primitive definition of architecture.

\begin{proposition}[Marked multiplicative factors give additive potential terms]\label{prop:additive-potential}
Suppose the positive raw coupling has an independently marked factorization
\[
K_{ri}(H)=\prod_{a=1}^{m}K_{ri}^{(a)}(H),
\qquad K_{ri}^{(a)}(H)>0.
\]
Then, with $\mathcal U_{ri}^{(a)}=-\tau\log K_{ri}^{(a)}$,
\[
-\tau\log K_{ri}=\sum_{a=1}^{m}\mathcal U_{ri}^{(a)}.
\]
\end{proposition}

\begin{proof}
Apply the logarithm to the product.
\end{proof}

Thus schema-, instance-, parameter-, and state-origin factors may contribute additive potential terms when those factors are independently marked. The extensional row does not uniquely determine such a provenance decomposition.

\subsection{Receiver free energy}\label{subsec:free-energy}

For $p\in\Delta(S_r)$, define
\begin{equation}\label{eq:receiver-free-energy}
\mathcal F_r[p;H]
:=-\sum_{i\in S_r}p_i s_{ri}(H)
+\tau D_{\mathrm{KL}}(p\Vert\pi_r).
\end{equation}
Equivalently,
\[
\mathcal F_r[p;H]
=
\sum_ip_i\mathcal U_{ri}(H)-\tau H(p).
\]

\begin{theorem}[Receiver free-energy characterization]\label{thm:receiver-free-energy}
Let
\[
\mathcal Z_r(H):=\sum_{i\in S_r}\pi_{ri}e^{s_{ri}(H)/\tau}.
\]
For every $p\in\Delta(S_r)$,
\begin{equation}\label{eq:free-energy-gap}
\boxed{
\mathcal F_r[p;H]
=
-\tau\log\mathcal Z_r(H)
+\tau D_{\mathrm{KL}}\!\left(p\Vert\alpha_r(\cdot\mid H)\right).
}
\end{equation}
Consequently, $\alpha_r(\cdot\mid H)$ is the unique minimizer of $\mathcal F_r[\cdot;H]$.
\end{theorem}

\begin{proof}
For admissible $i$,
\[
\log\frac{\alpha_{ri}(H)}{\pi_{ri}}
=\frac{s_{ri}(H)}{\tau}-\log\mathcal Z_r(H).
\]
Substitution into $D_{\mathrm{KL}}(p\Vert\alpha_r)$ and rearrangement gives \eqref{eq:free-energy-gap}. Nonnegativity of KL divergence gives the unique minimizer.
\end{proof}

The theorem characterizes the already-derived row. It does not postulate physical equilibrium for the network. It states that, conditional on the current receiver, admissible source sector, baseline, score, and row-simplex constraint, softmax is the unique allocation balancing evidence against departure from the reference measure. Other ensembles---balanced transport, column capacities, unnormalized flow, signed allocation, or hard minimum selection---would define different architectures over related score data \citep{cuturi2013sinkhorn}.

\subsection{Exact source-carrier quotienting}\label{subsec:coarse-potential}

An exact quotient of the source carrier preserves a full pushed-forward allocation and is different from the moment aggregation studied in \Cref{sec:interface-refinement}.

Let $q:S_r\twoheadrightarrow C_r$ be a source quotient. Define
\begin{equation}\label{eq:coarse-kernel}
K_r^{\mathrm{eff}}(c\mid H)
:=\sum_{q(i)=c}e^{-\mathcal U_{ri}(H)/\tau}
\end{equation}
and
\begin{equation}\label{eq:coarse-potential}
\mathcal U_r^{\mathrm{eff}}(c;H)
:=-\tau\log K_r^{\mathrm{eff}}(c\mid H).
\end{equation}

\begin{theorem}[Effective potential under a source quotient]\label{thm:potential-of-mean-force}
If $\bar\alpha_r=q_\#\alpha_r$, then
\[
\boxed{
\bar\alpha_r(c\mid H)
=
\frac{e^{-\mathcal U_r^{\mathrm{eff}}(c;H)/\tau}}
{\sum_{c'\in C_r}e^{-\mathcal U_r^{\mathrm{eff}}(c';H)/\tau}}.
}
\]
The partition sum is preserved.
\end{theorem}

\begin{proof}
The numerator is the sum of the microscopic Gibbs numerators over the fiber $q^{-1}(c)$. Summing over all quotient classes partitions the original source sum.
\end{proof}

The induced potential is a log-sum-exp over eliminated sources. If a quotient class contains $n_c$ sources with equal potential $u$, then
\[
\mathcal U_r^{\mathrm{eff}}(c)=u-\tau\log n_c.
\]
Microscopic multiplicity and potential combine. By contrast, ordinary value aggregation
\[
(\alpha_r,(v_i)_i)\longmapsto\sum_i\alpha_{ri}v_i
\]
generally retains only one selected moment and does not preserve the pushed-forward allocation.

\subsection{Gauges, observables, and heads}

\begin{proposition}[Receiver potential gauge]\label{prop:potential-gauge}
Two finite potentials on the same admissible row induce the same normalized allocation if and only if they differ by a receiver-dependent additive constant.
\end{proposition}

\begin{proof}
A common additive shift produces a common multiplicative Gibbs factor that cancels. Conversely, equality of normalized rows makes $e^{-\mathcal U_i'/\tau}/e^{-\mathcal U_i/\tau}$ independent of $i$.
\end{proof}

\begin{lemma}[Query--key coordinate gauge]\label{lem:qk-gauge}
For $R\in GL(d)$, the transformations
\[
q'(x)=Rq(x),
\qquad
k'(y)=R^{-\mathsf T}k(y)
\]
preserve every score $\ip{q(x)}{k(y)}$.
\end{lemma}

The potential gauge changes a landscape coordinate without changing the row. The query--key gauge changes coordinates without changing the score. Neither representation is uniquely individuated by routing behavior.

The value map attaches an observable to each source:
\[
z_r(H)=\E_{i\sim\alpha_r(\cdot\mid H)}[v_i(y_i(H))].
\]
Values are not generally part of the potential. Two heads may use the same allocation with different observables, or different allocations whose selected moments coincide.

\begin{corollary}[Multihead routed decomposition]\label{cor:multihead}
A multihead output
\[
z_r
=
\sum_{m\in\mathcal H}O_m\left(\sum_i\alpha_{ri}^{(m)}v_i^{(m)}\right)
\]
is a parallel family of head-indexed receiver-interface refinements followed by linear head recombination.
\end{corollary}

The heads do not form one joint probability law unless an additional joint normalization is imposed.

\subsection{Commitment ablation}

\begin{table}[t]
\centering
\small
\begin{tabularx}{\textwidth}{@{}lX@{}}
\toprule
Commitment changed & Neighboring receiver-interface regime \\
\midrule
Pairwise scalar evidence & higher-order, hyperedge, tuple, vector-, matrix-, or operator-valued contributions \\
Positive support & signed, inhibitory, or cancellation-based routing \\
Receiver-row anchoring & columnwise, global, balanced, capacity-constrained, sparse, hard, or unnormalized allocation \\
Source-local payload & pair-conditioned or receiver-conditioned messages \\
Additive aggregate & recurrent aggregation, nonlinear set state, or retained source-resolved memory \\
Shared finite separation rank & direct pair network, nonseparable kernel, or state-specific score coordinates \\
Fixed admissible sector & state-generated hard barriers or persistent learned relation states \\
One-shot exposure & iterative retrieval, recurrent source access, or global planning \\
\bottomrule
\end{tabularx}
\caption{The attention theorem as a map from architectural commitments to neighboring interface forms.}
\label{tab:attention-ablation}
\end{table}

The theorem therefore exposes, rather than eliminates, architecture. The Transformer endogenizes receiver-relative allocation but continues to fix the source carrier, pairwise scalar interaction language, positivity, row-simplex ensemble, value transport form, aggregate statistic, and one-step schedule.

\section{Local continuation and the feedforward sublayer}\label{sec:local-continuation}

Attention refines the receiver interface $Q_j$. The pointwise feedforward sublayer occupies the other side of the central factorization: it refines the receiver-local continuation $G_j$. This distinction explains both the exact hidden-coordinate computation and its architectural limit. The FFN can construct new local observables from the retained receiver state, but it cannot recover source distinctions already identified by $Q_j$.

\subsection{Marked refinements of a local continuation}

\begin{definition}[Marked local-continuation refinement]\label{def:local-continuation-refinement}
Let $G:Z\to W$ be a receiver-local continuation. A marked refinement of $G$ through a state space $A$ consists of maps
\[
Z\xrightarrow{a}A\xrightarrow{o}W
\]
with $G=o\circ a$, together with declared projections or coordinates on $A$ that are part of the architectural presentation.
\end{definition}

The mark is necessary. Every function can be padded with arbitrary identity maps, invertible recodings, or unused coordinates. A hidden state is architectural only when the specified continuation actually contains that intermediate presentation and its coordinate roles.

\subsection{Two-layer local continuation}

Consider the receiver-local map
\begin{equation}\label{eq:ffn}
\FFN(x)=W_2\phi(W_1x+b_1)+b_2,
\end{equation}
where $\phi$ acts coordinatewise. Let
\[
a(x):=\phi(W_1x+b_1)\in\R^{d_{\mathrm{ff}}}
\]
and let $E=\{1,\ldots,d_{\mathrm{ff}}\}$ index the marked activation coordinates. Write $w_e$ for the $e$th column of $W_2$.

\begin{proposition}[Marked hidden-coordinate decomposition]\label{prop:ffn-hidden-carrier}
The specified two-layer continuation factors through the marked activation presentation
\[
x\xrightarrow{a}\R^{d_{\mathrm{ff}}}\xrightarrow{y\mapsto W_2y+b_2}W,
\]
and
\begin{equation}\label{eq:ffn-direction-sum}
\boxed{\FFN(x)=b_2+\sum_{e\in E}a_e(x)w_e.}
\end{equation}
\end{proposition}

\begin{proof}
The coordinatewise activation produces the marked vector $a(x)=(a_e(x))_e$. Multiplication by $W_2$ forms the linear combination of its columns with those coefficients.
\end{proof}

The conclusion is stronger than a formal column expansion and weaker than a mechanism claim. The hidden coordinates are architectural because the declared two-layer sublayer computes and retains them as the intermediate input to $W_2$. The proposition does not establish that one hidden coordinate is an independent expert, causal mechanism, or uniquely individuated basis direction.

Unlike attention coefficients, the native $a_e(x)$ may be signed and need not have unit sum. A nonnegative normalized allocation can be manufactured by signed carrier doubling and amplitude extraction, but that is a derived representation rather than the native form. The exact construction is recorded in the supplement.

\subsection{Gated local continuation}

Many contemporary blocks use
\begin{equation}\label{eq:gated-ffn}
\GFFN(x)
=
W_o\left[\phi(W_gx+b_g)\odot(W_ux+b_u)\right]+b_o.
\end{equation}

\begin{proposition}[Gated hidden-coordinate decomposition]\label{prop:gated-ffn}
Let $w_e$ be the $e$th column of $W_o$ and define
\[
a_e(x)
:=
\phi((W_gx+b_g)_e)(W_ux+b_u)_e.
\]
Then
\[
\boxed{\GFFN(x)=b_o+\sum_e a_e(x)w_e.}
\]
The marked intermediate carrier is the coordinatewise gated activation vector.
\end{proposition}

\begin{proof}
The vector inside $W_o$ has coordinates $a_e(x)$, and $W_o$ forms the corresponding column combination.
\end{proof}

The gated form changes the local coefficient constructor while preserving the architectural pattern: receiver-local state constructs coefficients on a marked hidden coordinate family, then a shared output map recombines the resulting directions.

\subsection{Local observable construction and its limit}

Let the receiver interface after routed exposure be
\[
Q_i(H)\cong\bigl(h_i,z_i(H)\bigr),
\]
and let a residual and normalization map produce the FFN input $x_i$. Every pointwise FFN output is then a function of $x_i$. The sublayer can:
\begin{itemize}
\item construct nonlinear observables already determined by the retained receiver macrostate;
\item move those observables into coordinates that later restricted operations can use;
\item apply one shared local law across receiver occurrences, possibly conditioned on explicit type marks.
\end{itemize}
It cannot distinguish predecessor states that induce the same $x_i$. By \Cref{thm:squared-decomposition}, increasing local approximation capacity addresses restricted accessibility but not unresolved variation inside the interface fibers.

This yields the division of labor
\[
\boxed{
\begin{aligned}
\text{source-resolved routing}
&\text{ determines which nonlocal distinctions are constructed},\\
\text{aggregation and residual retention}
&\text{ determine which receiver state survives},\\
\text{the FFN}
&\text{ constructs local observables of that retained state}.
\end{aligned}}
\]

\subsection{Coordinate nonuniqueness is activation-relative}

For a linear factorization $W=AB$, every invertible intermediate change of basis $M$ gives $W=(AM^{-1})(MB)$. A nonlinear two-layer FFN does not admit arbitrary $GL(d_{\mathrm{ff}})$ changes of hidden basis, because a general $M$ does not commute with coordinatewise $\phi$.

The valid presentation symmetries are activation-compatible. They include hidden-unit permutations, with corresponding permutations of the rows of $W_1$ and columns of $W_2$. For positively homogeneous activations, suitable positive coordinate rescalings are also available when compensated between the incoming and outgoing weights. Further symmetries require separate proof for the chosen activation and parameterization. The paper therefore does not identify all algebraic factorizations of the external local map with the specified hidden architecture.

\subsection{A derived hidden potential is not foundational}

If one first doubles signed coordinates and normalizes their nonnegative magnitudes, the resulting distribution can be written in Gibbs form on its positive support by defining $-\tau_{\mathrm{hid}}\log p_x$. This is tautologically exact for any positive finite distribution. It is weaker than the attention result because the native FFN coefficients need not arise from a positive compositional evidence law, preserve one receiver-relative ratio structure, or place total amplitude inside the normalized row. The useful commonality is state-conditioned coefficient construction over a marked carrier, not a shared thermodynamic origin.

\section{The canonical Transformer block}\label{sec:full-block}

The preceding sections reconstructed the two sides of the architectural factorization separately. Attention constructs a receiver interface from the current token state. The residual, normalization, and FFN form a structured receiver-local continuation. This section states the whole block first at that coarse grain, then exposes a finer two-sublayer presentation.

\subsection{Canonical deterministic PreNorm block}

Let
\[
H=(h_i)_{i\in S}\in V^S
\]
be the current token state. Let $N_1,N_2:V\to V$ be receiver-local normalization maps and let $m\in\mathcal H$ index heads. Let $Q_m,K_m,V_m$, and $O_m$ be the declared linear head maps of compatible types. Define
\[
q_i^{(m)}=Q_mN_1(h_i),
\qquad
k_j^{(m)}=K_mN_1(h_j),
\qquad
v_j^{(m)}=V_mN_1(h_j).
\]
Let $\alpha_{ij}^{(m)}(H)$ be masked row-normalized attention coefficients. With optional attention bias $b_{\mathrm{att}}$, define
\begin{equation}\label{eq:token-route}
t_i(H)
:=
b_{\mathrm{att}}
+
\sum_{m\in\mathcal H}O_m\left(\sum_{j\in S}\alpha_{ij}^{(m)}(H)v_j^{(m)}\right).
\end{equation}
The first residual state is
\begin{equation}\label{eq:first-residual}
u_i:=h_i+t_i(H).
\end{equation}
For a two-layer local sublayer,
\begin{equation}\label{eq:local-ffn-block}
f_i(u_i):=W_2\phi(W_1N_2(u_i)+b_1)+b_2,
\end{equation}
and the block output is
\begin{equation}\label{eq:block-output}
h_i^+:=u_i+f_i(u_i).
\end{equation}

\subsection{Coarse receiver factorization}

\begin{theorem}[Receiver-interface and local-continuation decomposition]\label{thm:block-receiver-factorization}
For every receiver $i$, let
\[
Z_i^{\mathrm{att}}
:=
\im\bigl(H\mapsto(h_i,t_i(H))\bigr),
\]
and define the attention-constructed interface, corestricted to its realized image,
\begin{equation}\label{eq:block-coarse-interface}
Q_i^{\mathrm{att}}:V^S\twoheadrightarrow Z_i^{\mathrm{att}},
\qquad
Q_i^{\mathrm{att}}(H):=\bigl(h_i,t_i(H)\bigr).
\end{equation}
Define the local continuation on $Z_i^{\mathrm{att}}$ by
\begin{equation}\label{eq:block-coarse-continuation}
G_i^{\mathrm{block}}(h,t)
:=\nu+\FFN(N_2(\nu)),
\qquad
\nu:=h+t,
\end{equation}
where $\FFN(x)=W_2\phi(W_1x+b_1)+b_2$. Then the receiver branch of the canonical PreNorm block factors as
\[
\boxed{
B_i^{\mathrm{block}}
=
G_i^{\mathrm{block}}Q_i^{\mathrm{att}}.
}
\]
Moreover, $Q_i^{\mathrm{att}}$ has a source-resolved refinement through the head--source contributions in \eqref{eq:token-route}, while $G_i^{\mathrm{block}}$ has the marked hidden-coordinate refinement of \Cref{prop:ffn-hidden-carrier}.
\end{theorem}

\begin{proof}
Evaluating \eqref{eq:block-coarse-continuation} on \eqref{eq:block-coarse-interface} gives $\nu=h_i+t_i(H)=u_i$ and therefore
\[
G_i^{\mathrm{block}}Q_i^{\mathrm{att}}(H)
=u_i+W_2\phi(W_1N_2(u_i)+b_1)+b_2
=h_i^+.
\]
Linearity of each $O_m$ expands the transported exposure as the source-resolved family
\[
t_i(H)
=b_{\mathrm{att}}
+
\sum_{m,j}\alpha_{ij}^{(m)}(H)O_mv_j^{(m)}.
\]
The FFN refinement follows from the marked intermediate activation vector.
\end{proof}

The theorem gives the main architectural reading of the block:
\[
\boxed{
\text{attention constructs }Q_i
\qquad\text{and}\qquad
\text{the residual--FFN subprogram realizes }G_i.
}
\]
Biases remain affine terms. They need not be converted into artificial ``null sources'' merely to force every term into one routed sum.

\subsection{Fine two-sublayer presentation}

The coarse factorization treats the complete block as one receiver branch. A finer presentation marks the intermediate residual state $U=(u_i)_i$ as a successor state of the attention sublayer and predecessor state of the FFN sublayer.

\begin{corollary}[Fine two-stage receiver presentation]\label{cor:fine-two-stage-block}
The canonical block admits the finer composition
\[
H\xrightarrow{B^{\mathrm{att}}}U\xrightarrow{B^{\mathrm{ffn}}}H^+,
\]
where each displayed interface is corestricted to its realized image and
\[
Q_i^{\mathrm{tok}}(H)=\bigl(h_i,t_i(H)\bigr),
\qquad
G_i^{\mathrm{tok}}(h,t)=h+t,
\]
and
\[
Q_i^{\mathrm{ffn}}(U)=\bigl(u_i,f_i(u_i)\bigr),
\qquad
G_i^{\mathrm{res},2}(u,\delta)=u+\delta.
\]
The transported token contribution is source-resolved over head--source pairs. The second residual interface contains the completed local FFN contribution, while that contribution is internally refined by the actual hidden activation
\[
u_i
\longmapsto
a_i=\phi(W_1N_2(u_i)+b_1)
\longmapsto
f_i(u_i)=W_2a_i+b_2.
\]
Thus $Q_i^{\mathrm{ffn}}$ is not itself called a hidden interface: the hidden carrier occurs inside the construction of its second component.
\end{corollary}

\begin{proof}
The first factorization is \eqref{eq:first-residual}. The second is \eqref{eq:block-output}. Their composition equals the coarse branch in \Cref{thm:block-receiver-factorization}.
\end{proof}

The coarse and fine presentations answer different architectural questions. The coarse view identifies the information available to the final receiver continuation. The fine view exposes the intermediate state on which the second local sublayer operates. Neither grain is forced by the external function alone; the selected marks determine which factorization is being claimed.

\begin{corollary}[Gated local continuation]\label{cor:gated-block}
If \eqref{eq:local-ffn-block} is replaced by the gated map \eqref{eq:gated-ffn}, the coarse and fine block factorizations remain valid with hidden coefficients
\[
a_e(u_i)
=
\phi((W_gN_2(u_i)+b_g)_e)(W_uN_2(u_i)+b_u)_e.
\]
\end{corollary}

\subsection{Depth as repeated interface construction}

Let $B_\ell$ be the $\ell$th block and
\[
H^{(\ell+1)}=B_\ell(H^{(\ell)}).
\]
At each layer, the current state determines the evaluated interface constructors
\[
H^{(\ell)}
\longmapsto
\bigl(Q_{\ell,i}^{\mathrm{att}}(H^{(\ell)})\bigr)_i
\]
and, in the additive regime, a token-routing operator family
\[
H^{(\ell)}\longmapsto A_{\ell,H^{(\ell)}}.
\]
In the positive scalar specialization it also determines receiver- and head-relative potentials
\[
H^{(\ell)}\longmapsto\mathcal U_{\ell,H^{(\ell)}}.
\]
Ordinary function composition is enough to describe depth:
\[
H^{(0)}
\xrightarrow{B_0}
H^{(1)}
\xrightarrow{B_1}
\cdots
\xrightarrow{B_{L-1}}
H^{(L)}.
\]
No category of process cuts is required for the main reconstruction. Whether these represented transitions descend coherently from an implementation or remain equivalent under a change of cut is a separate question not assumed here.

\subsection{Conditionally variational routing, globally driven computation}

Within one positive attention row, the allocation $\alpha_{ri}^{(m)}$ uniquely minimizes the receiver free-energy functional of \Cref{thm:receiver-free-energy}. The whole block does not minimize one fixed global free energy. Its update changes the represented state used to construct the next scores, masks, values, interfaces, and local hidden coefficients.

The variational statement is therefore receiver- and state-conditional: each declared positive row minimizes its own current free-energy functional, while the block update changes the state from which later rows are constructed. No energy-conservation law, monotone global free-energy descent, or physical equilibrium interpretation follows.

\subsection{Residual continuation and schedule variants}

The residual path marks direct availability of the predecessor-local state:
\[
h_i\longmapsto h_i+t_i\longmapsto h_i+t_i+f_i.
\]
It prevents the routed aggregate from becoming the receiver's sole successor variable. This does not guarantee invertibility or information conservation; distinct predecessor states may still share one successor.

The exact $Q/G$ split depends on schedule. In a PostNorm block, normalization occurs after residual addition. In a parallel block, attention and FFN branches may receive the same predecessor state rather than a sequentially updated one. Reversed, recurrent, adaptive-depth, asynchronous, or carrier-changing schedules give different receiver-wise factorizations because they change which represented state is supplied to each interface constructor and continuation. The architecture is the marked schedule of these factorizations, not merely the multiset of formulas appearing in the block.

\section{Architectural coordinates and neighboring systems}\label{sec:coordinates}

The receiver-factorization view supports comparison without pretending that all architectures possess a scalar potential or one common message ontology. General coordinates belong to the $Q_j/G_j$ factorization. Potential, ensemble, and observable language is a specialization available only after positive scalar routing has been selected.

\subsection{Related descriptions and their primary objects}

\paragraph{Function classes and expressivity.}
Expressivity results ask which sequence maps or equivariant functions a model family can realize \citep{yun2020transformers,hu2025softmax}. The present paper fixes one represented update and asks which marked receiver obligations, interfaces, and local continuations realize it. Equal function classes therefore do not settle the architectural comparison developed here.

\paragraph{Programmatic and circuit descriptions.}
RASP and Transformer-circuit formalisms expose programs, algorithms, or internal computational motifs available in Transformer-like systems \citep{weiss2021thinking,elhage2021framework}. The receiver factorization is a prior architectural grammar: it distinguishes the receiver obligation, the information made available to it, the contribution structure before aggregation, and the continuation after aggregation. It does not replace a program- or circuit-level account of a particular trained system.

\paragraph{Relational, message-passing, and geometric descriptions.}
Graph-network and geometric-deep-learning frameworks foreground carriers, relational support, symmetry, and message laws \citep{battaglia2018relational,gilmer2017neural,bronstein2021geometric}. The present framework is compatible with those descriptions but separates four roles that a single message-passing formula may package together: source-resolved contribution, retained aggregate exposure, complete receiver interface, and receiver-local continuation.

\paragraph{Information availability and internal interpretation.}
Usable-information and conditional-probing frameworks ask which targets can be decoded under a selected computational class \citep{xu2020usable,hewitt2021conditional}. Here the interface fibers identify which predecessor distinctions are unavailable after a marked architectural map, while the restricted-locality term identifies what a selected continuation class fails to extract. Decodability alone does not establish that the probed interface is an implementation-warranted mechanism boundary.

\paragraph{Attention-specific descriptions.}
The attention result derived earlier is not a new layer formula and does not compete with the original definition of scaled dot-product attention \citep{vaswani2017attention}. It is a conditional normal form: the familiar row appears once the receiver interface is restricted to a declared package of scalar, positive, row-anchored, source-local, additive, and jointly separable commitments.

\subsection{General factorization coordinates}

\begin{longtable}{@{}p{0.18\textwidth}p{0.33\textwidth}p{0.42\textwidth}@{}}
\caption{General architectural coordinates of a receiver-factored update.}\label{tab:general-coordinates}\\
\toprule
Coordinate & Question & Examples of variation \\
\midrule
\endfirsthead
\toprule
Coordinate & Question & Examples of variation \\
\midrule
\endhead
Receiver presentation & Which successor obligations are separately addressed? & tokens, spatial sites, graph nodes, latent slots, experts, sectors, dynamically created carriers \\
Receiver interface $Q_j$ & Which predecessor distinctions reach receiver $j$? & full state, fixed neighborhood, recurrent summary, routed aggregate, source-resolved memory, multiple retained moments \\
Contribution language & How is transported information resolved before aggregation? & scalar weights, vector or operator contributions, pair-conditioned messages, higher-order tuples, recurrent access procedures \\
Aggregation & Which contribution distinctions survive into the receiver exposure? & sum, expectation, maximum, nonlinear set map, recurrent aggregation, exact quotient pushforward, no aggregation \\
Local continuation $G_j$ & What can the receiver construct from its interface? & affine map, residual addition, shared MLP, gated MLP, typed local program, stochastic local update \\
Composition and schedule & In what order are interfaces constructed and continuations applied? & sequential, parallel, reversed, recurrent, asynchronous, adaptive depth, carrier-changing continuation \\
\bottomrule
\end{longtable}

Provenance is an orthogonal annotation on every coordinate. A receiver relation, contribution operator, aggregation rule, local continuation, or schedule may be schema-fixed, instance-supplied, parameter-dependent, or state-conditioned. The marked source type and the derived essential-dependence set should not be conflated.

\subsection{Positive scalar specialization}

Inside the attention regime, the general coordinates specialize further:

\begin{longtable}{@{}p{0.20\textwidth}p{0.32\textwidth}p{0.40\textwidth}@{}}
\caption{Additional coordinates available in the positive scalar routing specialization.}\label{tab:positive-coordinates}\\
\toprule
Coordinate & Canonical Transformer choice & Neighboring possibilities \\
\midrule
\endfirsthead
\toprule
Coordinate & Canonical Transformer choice & Neighboring possibilities \\
\midrule
\endhead
Source configuration & supplied token or patch carrier with one admissible row per receiver and head & hierarchical, persistent, sparse, memory, latent-object, multiscale, or variable successor carriers \\
Score/potential family & pairwise scalar score with hard mask, positive baseline, and shared finite-rank query--key chart & direct pair networks, higher-order scores, nonseparable kernels, learned barriers, persistent relation states \\
Allocation ensemble & one positive simplex per receiver and head & column capacities, balanced or doubly stochastic transport, sparse or hard routing, unbalanced flow, globally coupled plans \\
Transported observables & source-local value vectors followed by linear head mixing & pair-conditioned messages, receiver-conditioned values, multiple observable families, operator-valued transport \\
Retained macrostate & one additive value expectation per head, then head recombination & full pushed-forward law, several moments, source provenance, recurrent receiver memory, adaptive-dimensional state \\
Score scale and schedule & fixed or parameterized scale inside a fixed block order & state-dependent scale, iterative allocation, reversed or parallel branches, conditional depth \\
\bottomrule
\end{longtable}

This table does not redefine architecture in thermodynamic terms. It states which additional design coordinates become visible once the general contribution language has been restricted to positive scalar allocation.

\subsection{Explicit realizations}

\paragraph{Convolution.}
Let $S$ be a finite lattice with a boundary convention and $D$ a finite offset set. A convolution
\[
(\operatorname{Conv}H)_r=\sum_{\delta\in D}W_\delta h_{r+\delta}
\]
has the additive routed realization
\[
C_{ri}
:=
\sum_{\substack{\delta\in D\\ i=r+\delta}}W_\delta,
\qquad
(\operatorname{Conv}H)_r=\sum_{i\in S}C_{ri}h_i.
\]
Under boundary conventions for which each offset labels a distinct source occurrence, this reduces to $C_{r,r+\delta}=W_\delta$. The summed definition also covers collisions created by periodic or quotient boundary conventions. The support and parameter tying are schema-fixed; the payload values are state-derived. If one instead keeps offset-labelled contributions as separate channels and forces a positive scalar interpretation, each such channel has singleton support, so no entropy-regularized competition within the channel remains. This realization clarifies the shared additive grammar without identifying convolution and attention as the same marked architecture \citep{cordonnier2020relationship}.

\paragraph{Graph attention.}
A graph-attention layer combines an instance-supplied admissible sector $N_\xi(r)$ with state-conditioned evidence inside that sector. Relative to full self-attention, it is less endogenous in admissibility but may admit richer edge types or message maps.

\paragraph{Kernel estimation.}
Normalized kernel estimation already realizes one positive scalar receiver row. Its contrast with a Transformer is the surrounding architecture: it normally computes one terminal estimate under a supplied comparison geometry rather than composing receiver-factored successor updates through depth.

\paragraph{Recurrent summaries.}
A recurrent model can realize long-range dependence through a low-dimensional persistent interface. It may use fewer source-resolved computations while creating stronger information barriers than an explicit attention row.

\paragraph{Mixture of experts.}
Expert routing constructs a token-to-expert interface. Capacity constraints may couple different tokens' assignments, so the allocation ensemble need not factor into independent receiver rows.

\paragraph{Operator-valued routing.}
An architecture may retain the general pair operators $C_{ri}(H)$ without separating scalar relevance from transport. Such a system remains within the receiver-interface framework but has no unique scalar potential or probability-row interpretation.

\subsection{What architecture comparison must preserve}

Function-class inclusion is too coarse. A comparison should declare which marks are being preserved, such as:
\begin{itemize}
\item receiver projections and carrier roles;
\item receiver interfaces and their refinements;
\item the distinction between source-resolved and aggregate state;
\item provenance typing and parameter-sharing rules;
\item local continuation structure;
\item schedule and successor presentation.
\end{itemize}
A universal specialization preorder would require additional decisions about harmless recoding, quotienting, carrier correspondence, and schedule comparison. Those decisions are claim-relative and require more comparison structure than the receiver factorization alone supplies. The present paper therefore uses explicit realization witnesses and coordinate profiles rather than asserting one total ordering of architectures.

The central comparison claim is:
\begin{quote}
The Transformer moves selected receiver-interface decisions into current-state computation while continuing to fix the receiver carrier, the admissible interaction language, the allocation ensemble, the transported observable family, the aggregate retained by each receiver, the local continuation family, and the block schedule.
\end{quote}

\section{Extension: state-conditioned successor presentations}\label{sec:variable-population}

The standard block constructs receiver-relative organization over a supplied carrier population. A further architectural move allows the marked successor presentation itself to vary with the current state. This changes which receiver obligations exist next rather than merely changing scores on one fixed receiver--source relation.

\subsection{Sector-indexed successor presentations}

Let $\mathfrak C$ be a finite or countable family of successor configurations. For each $c\in\mathfrak C$, let $R_c$ be its receiver set and let
\[
Y_c\subseteq\prod_{j\in R_c}W_{c,j}
\]
be the represented successor state in that sector. Define the tagged disjoint union
\begin{equation}\label{eq:sector-state-space}
Y^{\sqcup}:=\bigsqcup_{c\in\mathfrak C}Y_c
\end{equation}
and a deterministic update
\[
F:X\to Y^{\sqcup}.
\]
The sector projection
\[
\operatorname{sect}:Y^{\sqcup}\to\mathfrak C
\]
identifies the selected receiver presentation, while the point inside $Y_c$ gives the successor content.

\begin{definition}[Configurational endogenization]\label{def:configurational-endogenization}
The successor presentation is configurationally endogenous when the selected sector
\[
c_F:=\operatorname{sect}\circ F:X\to\mathfrak C
\]
is constructed from the current represented state rather than supplied independently of the update.
\end{definition}

This differs from relational endogenization on a fixed carrier. Standard self-attention changes receiver-relative allocation while preserving the ambient token population. Configurational endogenization permits the successor receiver set itself to merge, split, appear, disappear, change type, or change cardinality.

\subsection{Fixed-slot encoding does not settle the architecture}

\begin{proposition}[Finite-slot encoding]\label{prop:finite-slot-encoding}
Let $K<\infty$, let
\[
\mathcal X_{\le K}:=\bigsqcup_{N=0}^{K}V^N,
\]
and let $\bot\notin V$. The padding map
\[
\iota(v_1,\ldots,v_N)
=(v_1,\ldots,v_N,\bot,\ldots,\bot)
\in(V\sqcup\{\bot\})^K
\]
is injective on states whose null coordinates form a terminal suffix.
\end{proposition}

\begin{proof}
If a null coordinate occurs, its first position determines $N$; if none occurs, then $N=K$. The first $N$ entries then recover the unique variable-length state.
\end{proof}

The proposition shows that bounded variable-cardinality states can be represented extensionally by fixed slots. It does not determine whether occupancy is instance-supplied, generated by the current state, or merely used as a storage convention. Those systems can share one padded tensor space while differing in which successor obligations the architecture constructs and which inactive slots avoid later computation.

\subsection{Configuration and content are distinct coordinates}

For each $H\in X$, the tagged state $F(H)$ determines both
\[
c_F(H)\in\mathfrak C
\qquad\text{and}\qquad
F(H)\in Y_{c_F(H)}.
\]
The first coordinate selects the receiver presentation; the second supplies content inside that presentation. Identifying the sector label with the complete successor state would erase this distinction. For each $c\in\mathfrak C$, define the sector domain
\[
X_c:=c_F^{-1}(c).
\]
Whenever $X_c$ is nonempty, the restriction
\[
F_c:=F|_{X_c}:X_c\to Y_c
\]
is a fixed-sector represented update, and the receiver-factorization definitions of \Cref{sec:receiver-foundation} apply to its branch maps. Thus the variable-successor construction is a tagged family of ordinary receiver-factored updates together with the state-conditioned sector selector.

For multi-token generation, one may take
\[
Y_N=\Sigma^N,
\qquad
Y^{\sqcup}=\bigsqcup_{N=0}^{K}\Sigma^N.
\]
Then the selected count $N$ and the jointly constructed token block are separate architectural coordinates. This observation alone supplies no sampler, acceptance rule, equality with an autoregressive law, or compute improvement; it identifies the typed target required for a genuinely variable successor presentation. Stochastic realizations and their conditional Gibbs coordinates are collected in the supplement \citep{gloeckle2024multitoken,cai2024medusa,chen2023speculative}.

\section{Architectural presentation and stronger mechanism claims}\label{sec:mechanism-boundary}

The preceding sections construct exact architectural presentations. They do not by themselves establish that the marked interfaces are uniquely realized in an implementation, that a routing edge mediates an outcome under intervention, or that two presentations identify the same mechanism. This boundary specifies additional evidence obligations rather than weakening the architectural results.

\subsection{Function equality does not determine architecture}

Suppose a convolutional network, recurrent model, MLP, and Transformer realize the same external map $F:X\to Y$ on a declared domain. That equality does not determine:
\begin{itemize}
\item the marked receiver decomposition of $Y$;
\item the interfaces $Q_j$ through which the receiver branches factor;
\item whether source resolution is retained or discarded before local continuation;
\item whether the operative relation is supplied, parameter-static, or state-conditioned;
\item which local continuation family acts on the retained interface;
\item the schedule and successor presentation through depth.
\end{itemize}
Architecture comparison must therefore concern marked factorizations, not only realizable input--output functions. The re-presentation convention of \Cref{prop:interface-representation} removes harmless bijective coordinate changes of one interface, but it does not define a general equivalence relation over complete implementations.

\subsection{Routing coefficients do not by themselves establish mediation}

\begin{proposition}[Different routing rows can have the same aggregate]\label{prop:weights-not-mediation}
Let $v_1=v_2=v$, and let
\[
\alpha=\left(\frac34,\frac14\right),
\qquad
\alpha'=\left(\frac14,\frac34\right).
\]
Then $\alpha\neq\alpha'$ but
\[
\sum_i\alpha_iv_i
=
v
=
\sum_i\alpha_i'v_i.
\]
\end{proposition}

A routing coefficient belongs to the declared allocation rule. The proposition shows that changing that rule need not change even the immediate receiver aggregate. Consequently, a nonzero coefficient, low effective potential, or large score deformation does not alone establish that an isolated change in one source produces a receiver-level change under admissible controls. Equal payloads, aggregation collisions, or later continuation may remove the relevant variation.

\subsection{Additional evidence obligations}

Passing from an architectural presentation to an implementation-facing mechanism claim would require, at minimum:
\begin{itemize}
\item a declared implementation state space and implementation operations;
\item maps connecting implementation states and transitions to the represented update;
\item evidence that the marked receiver interfaces correspond to implemented occurrences or boundaries;
\item admissible interventions or contrasts and observable receiver outcomes;
\item invariance criteria distinguishing harmless re-presentation from constitutive change;
\item tracking rules when a carrier or role is claimed to persist through depth or across runs.
\end{itemize}
None of these follows from the existence of an exact factorization $B_j=G_jQ_j$. The present paper establishes the architecture conditional on its represented and marked data and blocks stronger inferences that those data alone do not support.

\section{Conclusion}\label{sec:conclusion}

The paper began after representation, with a map $F:X\to Y$. The missing architectural structure was not another formula but a marked receiver decomposition of the successor and a marked factorization of every receiver branch:
\[
\boxed{
B_j
=
\pi_j\chi_YF
=
G_jQ_j.
}
\]
Without the receiver marks, there are no separately addressed successor obligations. Without $Q_j$, receiver indexing with full shared access collapses to one product-valued map. Without $G_j$, the theory does not distinguish the information supplied to a receiver from the local computation performed after that information arrives. Bijective recodings of one interface are harmless re-presentations, while changes to its fibers alter which predecessor distinctions reach the receiver.

The factorization orders the later results. Interface fibers determine exact and approximate information barriers. Refining $Q_j$ into retained local state, source-resolved contributions, and aggregation distinguishes what is computed from what survives. In an additive linear regime, evaluated pair-contribution operators define a state-indexed operator family. Scalar relevance and transport become separate objects only when that factorization is itself marked.

Attention follows only after further restrictions. Pairwise scalar evidence, positive compositional linking, receiver-row ratio preservation, source-local payloads, additive aggregation, and shared finite separation rank yield masked query--key softmax. The supplement exhibits neighboring architectures obtained by removing these commitments one at a time. In the resulting positive scalar regime, effective-potential and free-energy expressions provide exact receiver-relative coordinates on an allocation that has already been derived; they are not the primitive definition of architecture.

The pointwise FFN enters on the other side of the central factorization. Its hidden coordinates are architectural because the specified local continuation computes a marked intermediate activation vector before output recombination. The FFN can construct nonlinear observables of the retained receiver state, but it cannot recover source distinctions removed before that state. A canonical PreNorm block is therefore, at a coarse grain, an attention-constructed receiver interface followed by a structured residual--FFN continuation. A finer grain exposes the intermediate residual state and the second receiver-wise update.

The resulting design space varies in receiver presentation, interface, contribution language, aggregation, local continuation, and schedule. Positive scalar routing adds a narrower set of allocation and potential coordinates. Variable successor presentations extend the same framework by allowing the marked target sector itself to depend on current state rather than treating a changed mask on one fixed carrier as creation of a new carrier.

Finally, an exact architectural presentation does not by itself establish an implementation boundary, an intervention-relative mediation relation, or identity across complete implementations. Those claims require additional evidence linking the marked presentation to implementation states, operations, admissible contrasts, invariance criteria, and tracking rules. The contribution here is the architectural object and the conditional reconstruction it supports: architecture comes before the familiar formula because the formula receives its architectural meaning from a marked organization of receiver access and continuation.


\bibliographystyle{plainnat}
\bibliography{references,references_additions}

\end{document}
